\section{Appendix}

%

\subsection{Higman Linearization of Algebraic Branching Programs}

In this section we present a proof of Theorem~\ref{block-higman}. More
precisely, we give an efficient deterministic algorithm to compute
Higman Linearization for a non-commutative polynomial given by an
algebraic branching program. We obtain this by generalizing the Higman
Linearization process (described in Theorem~\ref{ehigman}) to what we call
Block-Higman linearization. We restate Theorem~\ref{block-higman}) and
present its proof.

%
%
%
%

\begin{theorem}[Block Higman Linearization]\label{blockhigthm}
  Given as input an algebraic branching program of size $s$ computing
  a noncommutative polynomial $f\in\fR$, we can compute in
  deterministic polynomial time matrices
  $P, Q\in\fR^{\ell \times \ell}$ and a linear matrix
  $L\in\fR^{\ell\times \ell}$ such that
\begin{equation}\label{higeq-block-polynomial}
%
P\left(
\begin{array}{c|c}
f & 0 \\
\hline
0 & I_{\ell-1}
\end{array}
\right)Q ~=~L
\end{equation}
with $P$ upper triangular, $Q$ lower triangular, and the diagonal
entries of both $P$ and $Q$ are all $1$'s (hence, $P$ and $Q$ are both
units in $\fR^{\ell \times \ell}$). Futhermore, the algorithm computes
the entries of matrices $P$ and $Q$ as algebraic branching programs.
The entries of $P^{-1}$ and $Q^{-1}$ are also computable as algebraic
branching programs.

%
%
\end{theorem}

\begin{proof}
  Let $f\in \fR$ be the input noncommutative polynomial of degree $d$
  computed by an ABP of size $s$ and $d+1$ layers, where the $i^{th}$
  layer has say $n_i$ nodes. Then we have linear matrices
  $A_1,A_2,\ldots,A_d$ such that
  \[
    f=A_1\cdot A_2\cdots A_d,
  \]
  where $A_i$ is $n_i\times n_{i+1}$ for each $i$ and $n_1=n_{d+1}=1$.
  Note that the $(j,k)^{th}$ entry of $A_i$ is the linear form
  labeling the edge from $j^{th}$ node in layer $i$ to $k^{th}$
  node in layer $i+1$.

  We will prove the following more general result, dropping the
  constraint that $n_1=n_{d+1}=1$. Suppose $A_1,A_2,\ldots,A_d$ are
  linear matrices of compatible dimensions ($A_i$ is
  $n_i\times n_{i+1}$ for each $i$) such that the matrix product
\[
  M=A_1\cdot A_2\cdots A_d
\]  
is well defined. The algorithm we will describe will compute matrices
$P, Q, L$ such that

\begin{equation}\label{higeq-block-product}
%
P\left(
\begin{array}{c|c}
M & 0 \\
\hline
0 & I_{\ell-1}
\end{array}
\right)Q ~=~L
\end{equation}

and $P, Q, L$ have the properties as stated in the theorem.

%

%
%

The proof is by an easy induction on number $d$ of linear matrices
whose product is $M$. We set up this induction by describing a single
step of Block-Higman linearization writing the input matrix $M=AB$,
where $A=A_1\cdot A_2\cdots A_{d-1}$ and $B=A_d$ is a linear matrix.

So, let $M$ be the $r \times t$ matrix where $M=A\cdot B$ with
$A \in \fR^{r\times s}$ and $B \in \fR^{s\times t}$, and $B$ is a linear
matrix. We apply the following steps to $M$.

\begin{itemize}
\item Expand $M$ to a $(r+s) \times (t+s)$ matrix of the following
  shape by adding $s$ new rows and $s$ new columns, with the bottom
  right diagonal block being $I_s$ and the remaining entries zero to
  obtain the following:
\[
\left[
\begin{array}{c|c}
M & 0 \\
\hline
0 & I_s
\end{array}
\right]. 
\]
\item Use suitable block row and column operations to transform the
  matrix as follows
\[
\left(
\begin{array}{cc}
AB & 0 \\
0 & I_s
\end{array}
\right)\rightarrow
\left(
\begin{array}{cc}
AB & A \\
0 & I_s
\end{array}
\right)\rightarrow
\left(
\begin{array}{cc}
0 & A \\
-B & I_s
\end{array}
\right)
\]

Here the first step is realized by computing the matrix product
$A\cdot [0~|~I_s]$ and adding it to the respective blocks of the first
$r$ rows. The second step is realized by computing the matrix product
\[
 \left(
\begin{array}{c}
  A \\
  I_s
\end{array}
\right)\cdot (-B)
\]
and adding this to the respective blocks of the first $t$ columns.
These two steps are realized by left multiplication by $\left(
\begin{array}{c|c}
I_r & A \\
\hline
0 & I_s
\end{array}
\right)$ and right multiplication by $\left(
\begin{array}{c|c}
I_t & 0 \\
\hline
 -B & I_s
\end{array}
\right).$
\end{itemize}
So we have
\[\left(
\begin{array}{cc}
I_r & A \\
%
0 & I_s
\end{array}
\right)
\left(
\begin{array}{cc}
AB & 0 \\
0 & I_s
\end{array}
\right)
\left(
\begin{array}{cc}
I_t & 0 \\
%
 -B & I_s
\end{array}
\right) = 
\left(
\begin{array}{cc}
0 & A \\
-B & I_s
\end{array}
\right).
\]
In the above, we note that the row operation matrix $\left(
\begin{array}{cc}
I_r & A \\
%
0 & I_s
\end{array}
\right)$ is upper triangular with all diagonal entries $1$ where as the column operation matrix $\left(
\begin{array}{cc}
I_t & 0 \\
%
 -B & I_s
\end{array}
\right)$ is lower triangular with all diagonal entries $1$. 

Crucially, if $A=A_1\cdot A_2\cdots A_{d-1}$ and $B=A_d$, the
resulting matrix has only $A$ as the nonlinear block which is a
product of $d-1$ linear matrices. We can now apply induction to
the matrix $A$ to obtain the Block-Higman linearization of $M$
as claimed. 

We describe the intermediate steps of the induction in more detail, in
order to see the final shape of the linear matrix.

Let $M= A_1A_2\ldots A_d$ where $A_i \in \fR ^{n_i \times n_{i+1}}$,
are linear matrices. At the $i^{th}$ step of the induction
Block-Higman linearization transforms a matrix of the form $\left(
\begin{array}{cc}
* & A_1A_2\ldots A_{d-i} \\
%
* & *
\end{array} 
\right)$ into a matrix of the form $\left(
\begin{array}{cc}
* & A_1A_2\ldots A_{d-i-1} \\
%
* & *
\end{array} 
\right)$ 
where each $*$ indicates matrix blocks with linear entries. After
$d-1$ steps of Block-Higman linearization we obtain a linear matrix
which is an associate of $M$. 

In more detail, let $M_0=A_1A_2\ldots A_d, P_0=I_{n_1}$,
$Q_0=I_{n_{d+1}}$ and $t_0=0$. We have
$P_0 (M \oplus I_{t_0}) Q_0 = M_0$. Inductively, assume that we have
upper triangular matrix $P_i$ with all diagonal entries $1$ and a
lower triangular matrix $Q_i$ with all diagonal entries $1$ such that
\[ P_i \left(
\begin{array}{cc}
M & 0 \\
%
0 & I_{t_i}
\end{array} 
\right)Q_i = M_i.
\] 
Here, $M_i$ is a polynomial matrix which has top right block equal to
$A_1A_2 \ldots A_{d-i}$ and the other entries of $M$ are linear and
the entries of $P_i$ and $Q_i$ are all computable by ABPs. Let
$t_{i+1}=t_i + n_{d-i}$. Clearly,
$(P_i \oplus I_{n_{d-i}}) (M \oplus I_{t_{i+1}})(Q_i \oplus
I_{n_{d-i}})= M_i \oplus I_{n_{d-i}}$. Let $M_i'$ be a matrix obtained
from $M_i$ by replacing top right block by $0$. Let $M_{i+1}$ be a
$2\times 2$ block matrix with $M_i'$ as top left block and the
structure as shown below
\[
M_{i+1}= \left(
\begin{array}{ccccc|c|c}
 &&&&& 0& A_1A_2\ldots A_{d-i-1} \\
\hline
%
&&&&& & 0\\
&&&&& & \vdots\\
&&&&& & 0\\
%
\hline
0& 0& \ldots & 0 & 0& -A_{d-i} & I_{n_{d-i}}
\end{array} 
\right),
\]
where matrix blocks $-A_{d-i}$ and $A_1A_2\ldots A_{d-i-1}$ align with top right $0$ block in $M_i'$. Now we define suitable block row and column operations which transforms matrix $M_i \oplus I_{n_{d-i}}$ to $M_{i+1}$. 

By applying the Block-Higman linearization step we will obtain

%
%
%
%
%
%
%
%

%
%
%
%
%
%
%
%

\[
P'(M_i \oplus I_{n_{d-i}})Q' = M_{i+1},
\] 
where $P'$ and $Q'$ are upper and lower triangular matrices performing
the block row and column operations and their entries are computable
by ABPs. Letting $P_{i+1} = P' (P_i \oplus I_{n_{d-i}})$ and
$Q_{i+1}= (Q_i \oplus I_{n_{d-i}}) Q'$ we get
\[
P_{i+1} (M \oplus I_{t_{i+1}}) Q_{i+1} = M_{i+1}
\]
where $P_{i+1}$ and $Q_{i+1}$ are upper and lower triangular matrices
with diagonal entries $1$, the top right block (consisting of top
$n_1$ rows and last $n_{d-i}$ columns) of $M_{i+1}$ is
$A_1A_2 \ldots A_{d-i-1}$ and all other entries of $M_{i+1}$ are
linear. Continuing thus, we obtain upper and lower triangular matrices
$P_{d-1}$ and $Q_{d-1}$ with all diagonal entries $1$ such that
$P_{d-1} (M \oplus I_{t_{d-1}})Q_{d-1}= M_{d-1}$ which is a linear
matrix. Moreover, it is easy to see that entries of $P_{d-1}$ and
$Q_{d-1}$ are given by polynomial size ABPs (of $O(s^2)$ size to be
precise).

Carefully observing the shapes of the matrices $M_i$ we note that the
final linearized matrix $M_{d-1}$ has the form
\[
\left(
\begin{array}{ccccccc}
0 & 0 & 0 & 0 & \ldots & 0 & A_1\\
%
A_d & I_{n_{d}} & 0 & 0 & \ldots & 0 & 0\\
0 & A_{d-1} & I_{n_{d-1}} & 0 & \ldots & 0 & 0 \\
0 & 0 & A_{d-2} & I_{n_{d-2}}& \ldots & 0 & 0 \\
0 & 0 & \ldots & \vdots & \ddots & \vdots & 0\\
0& 0 & 0& 0 & \ldots & A_2 & I_{n_2}
\end{array} 
\right)
\]

$M_{d-1}$ is a $d\times d$ block matrix with
\begin{itemize}
\item[-]  $n_1\times n_2$ sized top right block in $A_1$.
\item[-]  $(i,i)^{th}$ block is $I_{n_{d+2-i}}$ for $2\leq i \leq d$.
\item[-] $(i,i-1)^{th}$ block is $A_{d+2-i}$ for $2\leq i \leq d$.
\item[-] all other entries are $0$.
\end{itemize}

\end{proof}


\subsection{Missing proofs from Section~\ref{prelim}}

\begin{proofof}{Theorem~\ref{cohnthm}}
Let $C\in\fR^{d\times d}$ be a full and right monic linear
matrix. Suppose Equation~\ref{eq-cohnthm} holds for some invertible
scalar matrices $S,S'$. Then we can write

\[
SCS'=  \left(
\begin{array}{cc}
A & 0 \\
D & B
\end{array}
\right) =   \left(
\begin{array}{cc}
A & 0 \\
0 & I
\end{array}
\right)\cdot 
  \left(
\begin{array}{cc}
I & 0 \\
D & I
\end{array}
\right)\cdot
\left(
\begin{array}{cc}
I & 0 \\
0 & B
\end{array}
\right).
\]
Since $C$ is right monic and $S, S'$ are invertible scalar matrices
the linear matrix $SCS'=\left(\begin{array}{cc} A & 0 \\ D & B
\end{array}\right)$ is also full and right monic. Writing it as
\[
\left(
\begin{array}{cc}
A & 0 \\
D & B
\end{array}
\right) =   \left(
\begin{array}{cc}
A_0 & 0 \\
D_0 & B_0
\end{array}
\right) + \sum_{i=1}^n \left(
\begin{array}{cc}
A_i & 0 \\
D_i & B_i
\end{array}
\right)\cdot x_i,
\]
it means the matrix
\[
\left[
\begin{array}{cc|}
A_1 & 0 \\
D_1 & B_1
\end{array}
\begin{array}{cc|}
A_2 & 0 \\
D_2 & B_2
\end{array}
 \ldots
\begin{array}{|cc}
A_n & 0 \\
D_n & B_n
\end{array}
\right]
\]
is full row rank. With suitable row operations applied to the above we
can see that both $[A_1 A_2\ldots A_n]$ and $[B_1 B_2 \ldots B_n]$ are
full row rank. Therefore, both $A$ and $B$ are full right monic
matrices hence they are nonunits by Lemma~\ref{monic-nonunit}.  Hence $A\oplus I$ and $B\oplus I$ are both non-units  which
implies that the factorization of $SCS'$ is nontrivial and hence $C$
is not an atom. 

Conversely, suppose $C$ is not an atom and $C=F\cdot G$ is a
nontrivial factorization. That means both $F$ and $G$ are full and
non-units. As $C$ is a linear matrix, applying \cite[Lemma
  5.8.7]{Cohnfir} we can assume that both $F$ and $G$ are linear
matrices. Now, since $F$ is a full linear matrix, by
Theorem~\ref{full-monic} (and Remark~\ref{remark-full-monic}) there
are a scalar invertible matrix $S_1$ and polynomial matrix $U_1$,
which is a unit, such that $S_1FU_1 = A\oplus I$ such that $A$ is
\emph{left} monic. Therefore, we have
\[
S_1 C = S_1 F U U^{-1} G = \left(
\begin{array}{cc}
A & 0 \\
0 & I
\end{array}
\right) \cdot \left(
\begin{array}{cc}
G'_1 & G'_3 \\
G'_2 & G'_4
\end{array}
\right) = \left(
\begin{array}{cc}
AG'_1 & AG'_3 \\
G'_2 & G'_4
\end{array}
\right).
\]
As $S_1 C$ is a linear matrix and $A$ is a left monic linear matrix we
can assume that $G'_1$ and $G'_3$ are scalar matrices. Since $S_1 C$ is
full rank, it forces the matrix $[G'_1 G'_3]$ to be full row rank (say
$r$, where $A$ is $r\times r$). Therefore, there is an invertible
scalar matrix $S'$ such that $[G'_1 G'_3]S' = [I_r 0]$. Putting it together,
we get the factorization
\[
SCS' = \left(
\begin{array}{cc}
A & 0 \\
0 & I
\end{array}
\right) \cdot
\left(
\begin{array}{cc}
I_r & 0 \\
G''_2 & G''_4
\end{array}
\right) =  \left(
\begin{array}{cc}
A & 0 \\
G''_2 & G''_4
\end{array}
\right)
\]
as claimed by the theorem.
\end{proofof}



\subsection{Proof of Lemma~\ref{hkv20-main-lemma}}


We present a self-contained proof of Lemma~\ref{hkv20-main-lemma} of
\cite{HKV20}.

\begin{definition}\label{def-setminus}
Let $U, V \subseteq \F^D$ be subspaces of $\F^D$ and $d= \dim U$. Fix
a basis $u_1, u_2, \ldots, u_\ell \in \F^D$ for $U \cap V$ and extend
it to a basis $u_1, u_2, \ldots, u_\ell, u_{\ell+1}, \ldots, u_d$ for
$U$. Further, let $u_1, u_2, ..., u_D$ be a basis for $\F^D$ obtained
by extending the above basis for $U$. Then $U\setminus V$ is defined
as $span(u_{\ell+1}, u_{\ell+2}, \ldots, u_d)$, i.e.
\[
U\setminus V = \{ \sum_{i= \ell+1}^{d} \alpha_i u_i | \alpha_i \in \mathbb{F} \textrm{ for } \ell < i \leq d \} 
\]
\end{definition}

Clearly $\dim U\setminus V = \dim U - \dim U \cap V$. Notice that
although the subspace $U\setminus V$ is basis dependent, the number
$\dim U\setminus V$ is independent of the construction of $U\setminus
V$.

\begin{definition}\label{projection}
Let $\mathcal{U}=\{U_1, U_2, \ldots, U_d \}$ be a collection of
subspaces of $\mathbb{F}^{D}$. For each $i \in [d]$ define
$\hat{U_i}^{(\mathcal{U})} = U_i \setminus (\sum_{k \neq i} U_k)$ as
above with respect to fixed bases for the subspaces.
\end{definition}

We first prove a technical lemma, essentially using the
inclusion-exclusion principle.

\begin{lemma}\label{hkv-rank-lemma}
Let $\mathcal{U}= \{U_1, U_2, \ldots, U_d\}$ be a collection of
subspaces of $\mathbb{F}^D$ for $d\geq 1$. Then
\[
\sum_{i=1}^{d}~ \left[\dim U_i + \dim \hat{U}_i^{(\mathcal{U})}
  \right]~~ \geq ~~2~\cdot~\dim \sum_{i=1}^d U_i.
\]
\end{lemma}

\begin{proof}
The proof will be by induction on $d$. The base case, $d=1$, is
obvious. Suppose it is is true for all $t <d$. I.e. for any subspace
collection $\mathcal{V} = \{V_1, V_2, \ldots, V_t\}$ we have
\[
\sum_{i=1}^{t}~ \left[\dim V_i + \dim  \hat{V}_i^{(\mathcal{V})} \right]~~ \geq ~~2~\cdot ~\dim\sum_{i=1}^t V_i.
\]
Letting $V_i= U_i$ for $1\leq i \leq d-2$ and $V_{d-1}= U_{d-1}+U_d$
in the above, we have
\[
\sum_{i=1}^{d-1}~ \left[\dim V_i + \dim \hat{V}_i^{(\mathcal{V})}
  \right] ~~\geq ~~2~\cdot~ \dim\sum_{i=1}^{d-1} V_i = 2\cdot\dim\sum_{i=1}^d U_i. 
\]

For the induction we need to show that $\sum_{i=1}^{d-1} (\dim V_i +
\dim \hat{V}_i^{(\mathcal{V})})\leq \sum_{i=1}^{d} (\dim U_i + \dim
\hat{U}_i^{(\mathcal{U})})$.

Now, \[\sum_{i=1}^{d-1}~ (\dim V_i + \dim
  \hat{V}_i^{(\mathcal{V})})\] is
\begin{eqnarray*}
 &=& \dim V_{d-1} + \dim \hat{V}_{d-1}^{(\mathcal{V})} + \sum_{i=1}^{d-2}~ \left[  ~  \dim U_i + \dim(U_i \setminus ( U_{d-1}+U_{d} + \sum_{k \neq i, k<d-1}U_k ))~ \right]  \\
 &=&  \dim V_{d-1} + \dim \hat{V}_{d-1}^{(\mathcal{V})} + \sum_{i=1}^{d-2}~ \left[ ~\dim U_i + \dim( U_i \setminus \sum_{k \neq i, k\leq d} U_k) \right] \\
 &=&  \dim V_{d-1} + \dim \hat{V}_{d-1}^{(\mathcal{V})}+ \sum_{i=1}^{d-2} ~\left[ \dim U_i + \dim \hat{U}_i^{(\mathcal{U})} \right]\\
 &=& \dim U_{d-1} + \dim U_d -dim( U_{d-1} \cap U_d ) + \dim \hat{V}_{d-1}^{(\mathcal{V})}+ \sum_{i=1}^{d-2}~\left [ \dim U_i + \dim \hat{U}_i^{(\mathcal{U})} \right] \\
  &=& \left[ \sum_{i=1}^{d} \dim U_i \right] + \left[ \sum_{i=1}^{d-2}~ \dim \hat{U}_i^{(\mathcal{U})} \right] + \dim \hat{V}_{d-1}^{(\mathcal{V})}- \dim (U_{d-1} \cap U_d).
\end{eqnarray*}

%
%
%

Hence, to complete the proof it suffices to show the following claim.

\begin{claim}
\[ \dim \hat{V}_{d-1}^{(\mathcal{V})}~ \leq~ \dim  \hat{U}_{d-1}^{(\mathcal{U})}  + \dim  \hat{U}_d^{(\mathcal{U})} + \dim (U_{d-1} \cap U_d) \]
\end{claim}
\begin{proofof}{Claim}
Let $T= U_1 + U_2 + \ldots + U_{d-2}$. Let $D, D_1, D_2$, and $D_3$ denote
dimensions of $T+U_{d-1}+U_d, U_{d-1}, U_d$, and $T$ respectively.

We have 
\begin{eqnarray*}
\dim \hat{V}_{d-1}^{(\mathcal{V})} &=& \dim (U_{d-1}+U_d)-\dim ((U_{d-1}+U_d) \cap T)\\
&=&\dim (U_{d-1}+U_d)-\dim (U_{d-1}+U_d)-\dim (T) + \dim (U_{d-1}+U_d+T)\\
&=& D-D_3
\end{eqnarray*}

%
%
%
%

Similarly,
\begin{eqnarray*}
\dim  \hat{U}_{d-1}^{(\mathcal{U})}  &=& \dim (U_{d-1}) - \dim (U_{d-1} \cap (T+ U_d))\\
&=&D_1 + \dim (T+U_{d-1}+U_d)- \dim (U_{d-1}) - \dim (T+U_d)\\
&=&D_1 +D -D_1- \dim (T)-\dim (U_d)+\dim (T \cap U_d)\\
&=& D-D_3-D_2 + \dim(T \cap U_d). 
\end{eqnarray*}
Likewise, we also have
\[
\dim ( \hat{U}_{d}^{(\mathcal{U})} )=D-D_3-D_1+ \dim (T \cap U_{d-1}).
\]
It is clear that the claim is equivalent to
\[
 D \geq D_1 + D_2 + D_3 - \dim (T \cap U_d) - \dim (T \cap U_{d-1})-
 \dim (U_{d-1} \cap U_d)
 \]
which follows immediately from the Inclusion-Exclusion Principle. 
\end{proofof}
\end{proof}

%
%
%

Now we will present a complete proof for Lemma \ref{hkv20-main-lemma}
from \cite{HKV20} where the proof is sketchy.\smallskip
 
{\bf Proof of Lemma \ref{hkv20-main-lemma} }

Let $L=A_0 + \sum_{i=1}^n A_i x_i$ Where $A_i \in \mathbb{F}_q^{d
  \times d}$ for $i \in [n]$. So, $L'= A_0 \otimes I_\ell +
\sum_{i=1}^{n} A_i \otimes Y_i$ . From standard properties of the
Kronecker product of matrices, there are a row permutation matrix $R$
and a column permutation matrix $C$ such that $L''= RL'C = I_\ell
\otimes A_0 + \sum_{i=1}^n Y_i \otimes A_i$. So $L'' = I_\ell \otimes
A_0 + \sum_{i=1}^n \sum_{1 \leq j, k \leq \ell} (E_{j,k} \otimes A_i)~
y_{i,j,k}$, where $E_{j,k}$ is a a matrix with $(j,k)^{th}$ entry one
and rest all entries equal to zero.
     
 We have $G L' H = \left(
\begin{array}{c|c}
A' & 0 \\
\hline
D' & B'
\end{array}
\right)$
Where $A'$ is $d' \times d'$ linear matrix for $d'>0$, $B'$ is $d''
\times d''$ linear matrix with $d' + d'' = d \ell$. Hence, $GR^{-1}
L''C^{-1} H = \left(
\begin{array}{c|c}
A' & 0 \\
\hline
D' & B'
\end{array}
\right)$.  Let $GR^{-1}=P_0$ and $C^{-1}H=Q_0$. Let $[ P_1 P_2 \ldots
  P_\ell]$ be the (full row rank) matrix obtained by picking the top
$d'$ rows of $P_0$ where each $P_i$ is $d' \times d$ scalar
matrix. Similarly let $[ Q_1^T Q_2^T \ldots Q_\ell^T]^T$ be the (full column
rank) matrix obtained by picking the rightmost $d''$ columns of $Q_0$
where each $Q_i$ is $d \times d''$ scalar matrix. Clearly,
\[   [ P_1 P_2 \ldots P_\ell] L'' [ Q_1^T Q_2^T \ldots Q_\ell^T]^T = 0, \]
which implies,
\[   [ P_1 P_2 \ldots P_\ell]\left[ I_\ell \otimes A_0 + \sum_{i=1}^n \sum_{1 \leq j, k \leq \ell} (E_{j,k} \otimes A_i)~ y_{i,j,k} \right] [ Q_1^T Q_2^T \ldots Q_\ell^T]^T = 0. \]

Equating the coefficients of each $y_{i,j,k}$ to zero we get the following.

\begin{equation} \label{eqn-A0}
\sum _{i=1}^\ell P_i A_0 Q_i = 0.
\end{equation} 

\begin{equation} \label{eqn-An0}
P_j A_i Q_k =0 \textrm{ for each } i>0 \textrm { and } 1\leq j,k \leq
\ell.
\end{equation} 

For each $i\in [\ell]$ the matrix $P_i$ is a linear transformation
from $\F^d$ to $\F^{d'}$. Let $U_i= \range(P_i)=\{P_i u | u \in
\F^{d}\}$ for each $i$, and $\mathcal{U}= \{U_1, U_2, \ldots, U_{\ell}
\}$. Let $T_i = \sum _{j \neq i} U_j$ for $i\in [\ell]$. Clearly,
$U_1+ U_2 + \ldots + U_\ell= \mathbb{F}_q^{d'}$ as $[ P_1 P_2\ldots
  P_\ell]$ is a full row rank matrix. For $i\in[\ell]$, let
$\hat{P}_i$ be a linear transformation from $\mathbb{F}^{d'}$ to
$\mathbb{F}^{d'}$ defined as follows. Fix a basis $u_{i,1}, u_{i,2},
\ldots, u_{i,r_i}$ of the subspace $U_i \cap T_i$. Extend it to a
basis $u_{i,1}, u_{i,2}, \ldots, u_{i,r_i}, u_{i,r_i+1}, \ldots,
u_{i,k_i}$, $k_i \geq r$, for $U_i$. Further, extend this basis of
$U_i$ to a complete basis $u_{i,1}, u_{i,2}, \ldots, u_{i,d'}$, for
$\mathbb{F}^{d'}$, where $d' \geq k_i$. For any vector $u =
\sum_{j=1}^{d'} \alpha_j u_{i,j}$ in $\mathbb{F}^{d'}$ let $\hat{P}_i
(u) = \sum_{j=r+1}^k \alpha_j u_{i,j}$. So, $\hat{P}_i(u)$ is the
vector obtained by projecting to the subspace $U_i\setminus T_i$
(which is defined w.r.t.\ the above basis). Hence,
$\hat{P}_i(u_{i,t})= u_{i,t}$ for $r_i < t \leq k_i$ and
$\hat{P}_i(u_{i,t})= 0$ otherwise. This defines a $d' \times d'$
matrix for each $\hat{P}_i$ for $i \in [\ell]$, which we also refer to
as $\hat{P}_i$ by abuse of notation. From the Definition
\ref{projection}, it follows that $\range(\hat{P}_i) =
\hat{U}_i^{(\mathcal{U})}$, so $\rank(P_i)=\dim
\hat{U}_i^{(\mathcal{U})}$.  Clearly, $\rank(\hat{P}_i P_i) =
\rank(\hat{P}_i)$ for $i \in [\ell]$. Now, by Lemma
\ref{hkv-rank-lemma} applied to the collection $\mathcal{U}= \{U_1,
U_2, \ldots, U_{\ell} \}$ we get
\[
\sum_{i=1}^{\ell}~ \left[ \rank P_i ~+~ \rank \hat{P}_i \right]~~ \geq
~~ 2\cdot ~ \dim \sum_{i=1}^{\ell} \range(P_i) = 2d'.
\]      

Similarly, each $Q_i:\F^{d''}\to \F^d$ is a linear map. We can define
the corresponding linear maps $\hat{Q}_i : \mathbb{F}^{d''} \to
\mathbb{F}^{d''}$ and associated $d'' \times d''$ sized matrices and
we will have $\rank Q_i\hat{Q}_i = \rank \hat{Q_i}$ for each
$i$. Applying the above argument we will get
 \[
\sum_{i=1}^{\ell}~ \left[ \rank Q_i ~+~ \rank \hat{Q}_i \right]~~ \geq
~~ 2\cdot ~ \dim \sum_{i=1}^{\ell} \range(Q_i) = 2d''.
\] 

Adding the two inequalities yields

\begin{equation}\label{eqn-php}
\sum_{i=1}^{\ell}~ \left( \rank \hat{P}_i ~+~ \rank Q_i \right)~ +  \left( \rank {P}_i ~+~ \rank \hat{Q}_i \right) ~~\geq ~~ 2\cdot ~ (d'+d'')= 2d\ell.
\end{equation}
From the Equation \ref{eqn-php} we would like to prove the following Claim.

\begin{claim} \label{php-simplify}
There exist index $i \in[\ell]$ such that $\rank \hat{P}_i ~+~ \rank Q_i \geq d$ and $\rank \hat{P}_i$, $\rank Q_i>0$ or $\rank {P}_i ~+~ \rank \hat{Q}_i \geq d$ and $\rank {P}_i$, $\rank \hat{Q}_i>0$. 
\end{claim}
First we complete the proof of the Lemma \ref{hkv20-main-lemma} assuming the Claim \ref{php-simplify}.
Without loss of generality, let index $i=1$ satisfies the Claim \ref{php-simplify} and further, let $\rank
\hat{P}_1 + \rank Q_1 \geq d$ with $\rank \hat{P}_1, \rank Q_1 >0$
(other case handled similarly).

Equation \ref{eqn-A0} implies that $\range(P_1 A_0 Q_1) \subseteq
T_1$, also clearly, $\range(P_1 A_0 Q_1) \subseteq \range(P_1)=
U_1$. Which implies $\range(P_1 A_0 Q_1)\subseteq U_1 \cap T_1$. Hence
$\hat{P}_1 P_1 A_0 Q_1 =0$. Equation \ref{eqn-An0} implies that,
$\hat{P}_1 P_j A_i Q_k= 0$ for all $i \geq 1$ and $1\leq j, k \leq
\ell$. So we get $\hat{P}_1 P_1 L Q_1 = 0$. Now $\rank(\hat{P}_1P_1) =
\rank(\hat{P}_1)\geq 1$, $\rank(Q_1) \geq 1$ and $\rank(\hat{P}_1
P_1)+\rank(Q_1) \geq d$. It follows that there exist $0 < e' \leq
\rank(\hat{P}_1 P_1)$ and $0 < e'' \leq \rank(Q_1)$ with $e'+e'' =
d$. By choosing $e'$ linearly independent rows of $\hat{P}_1P_1$ and
$e''$ linearly independent columns of $Q_1$ we obtain full row rank
matrix $U' \in \mathbb{F}_q^{e' \times d}$ and a full column rank
matrix $V' \in \mathbb{F}_q^{d \times e''}$ respectively.  Now we
extend $U'$ to a $d \times d$ matrix $U$ by adding any $d-e'$ linearly
independent rows such that $U$ is invertible. Similarly we extend $V'$
to a $d \times d$ matrix $V$ by adding any $d-e''$ linearly
independent columns such that $V$ is invertible. We clearly have $ULV
=\left(
\begin{array}{c|c}
A & 0 \\
\hline
D & B
\end{array}
\right)$ for some linear matrices $A, D, B$ such that $A \in \fR^{e' \times e'}$ and $B \in \fR^{e'' \times e''}$ with $0< e', e''$ and $e' + e'' =d$ as required. This completes the proof of Lemma \ref{hkv20-main-lemma}.
\\

\begin{proofof}{Claim~\ref{php-simplify}}

Each $P_i$ has $d$ columns and each $Q_i$ has $d$ rows. Thus, $\rank
P_i\le d$ and $\rank Q_i\le d$. Also, $\rank \hat{P}_i \leq \rank P_i$
and $\rank \hat{Q}_i \leq \rank Q_i$. Hence, $\rank \hat{P}_i + \rank
Q_i\le 2d$ and $\rank P_i + \rank \hat{Q}_i\le 2d$ for each $i$.

It follows from Inequality \ref{eqn-php} that if there is an $i$ for
which either $\rank \hat{P}_i + \rank Q_i <d$ or $\rank {P}_i + \rank
\hat{Q}_i< d$ then there must be an index $j$ such that either $\rank
\hat{P}_j + \rank Q_j > d$ or $\rank {P}_j + \rank \hat{Q}_j > d$. Two
cases arise:

\begin{enumerate}
\item for all $j\in [\ell]$, $\rank \hat{P}_j + \rank Q_j = d$ and
  $\rank {P}_j + \rank \hat{Q}_j = d$.
\item there is $j \in [\ell]$ with either $\rank \hat{P}_j + \rank Q_j
  > d$ or $\rank {P}_j + \rank \hat{Q}_j > d$.
\end{enumerate}

Suppose the first case occurs. It has the following two subcases.

\begin{enumerate}
\item[(a)] for all $j \in[\ell]$, $\rank \hat{P}_j=0$ or $\rank Q_j = 0$ and $\rank {P}_j=0$ or $\rank \hat{Q}_j=0$.
\item[(b)] there is $j\in [\ell]$ such that $\rank \hat{P}_j, \rank
  Q_j >0$ or $\rank {P}_j, \rank \hat{Q}_j> 0$
\end{enumerate}

First, consider Case 1(a). Note that $\rank \hat{P}_j=0$ implies
$\rank Q_j =d$. And $\rank Q_j=0$ implies $\rank \hat{P}_j=d$, which
implies $\rank P_j =d$. Thus, either $\rank P_j =d$ or $\rank Q_j =d$
for every $j$. Moreover, Case 1(a) also implies $\rank P_j, \rank
Q_j\in \{0,d\}$ for each $j$.

Now as $[P_1 P_2 \ldots P_\ell]$ has full row rank and $[Q_1^T
  Q_2^T~\ldots Q_\ell^T]^T$ has full column rank, Case 1(a) implies
that there are indices $j,k \in [\ell]$ such that $P_j$ and $Q_k$ are
both rank $d$ matrices. As $P_j$ is full column rank matrix, there is
a $d \times d'$ matrix $P_j'$ such that $P_j' P_j = I_d$. Similarly,
there is a $d'' \times d$ matrix $Q_k'$ such that $Q_k Q_k'=I_d$. Now
from Equation \ref{eqn-An0} we know that $P_j A_i Q_k = 0$ for all
$i$, $1\leq i \leq n$. Hence $P_j' P_j A_i Q_k Q_k' = 0$ for all $i$,
$1\leq i \leq n$. Consequently, $A_i=0$ for $1\le i\le n$ which is a
contradiction to the lemma statement. Hence case 1(a) cannot occur.

If case 1(b) or 2 holds then for some index $j\in [\ell]$ either
$\rank \hat{P}_j + \rank Q_j \geq d$ with $\rank \hat{P}_j$, $\rank
Q_j >0$ or $\rank {P}_j + \rank \hat{Q}_j \geq d$ with $\rank {P}_j$,
$\rank \hat{Q}_j>0$. 
\end{proofof}
